% GENERATED FILE. Edit canonical lessons, then run npm run generate:books.
% canonical-source-hash: fnv1a64:4a0318876a5db498
% canonical-lessons: KA-C400-bajet, KA-C400-adaya, KA-C400-sabsidi, KA-C400-viniyogisu, KA-C400-vyartha-madu

\chapter{Public Money}
\label{ch:ka-public-money}
\hlchaptermodality{400}

\begin{chapteropening}
\textbf{By the end of this chapter:} \emph{I can say a budget, income, a subsidy, to utilise and to waste.}

Complete the last lesson of chapter 400: I can say a budget, income, a subsidy, to utilise and to waste.
\end{chapteropening}

\section[\texorpdfstring{\kn{ಬಜೆಟ್} (bajeṭ)}{ಬಜೆಟ್ (bajeṭ)}]{\kn{ಬಜೆಟ್} (bajeṭ) — a budget}

\label{lesson:KA-C400-bajet}



\begin{quote}
Before the new one: say the Kannada for to distribute, then the Kannada for rural.
\end{quote}

\subsection*{You'll want to know: \kn{ಬಜೆಟ್}}
\textbf{\kn{ಬಜೆಟ್}} — \emph{bajeṭ} — \textquotedblleft{}a budget\textquotedblright{}.

\begin{quote}
\kn{ಈ} \kn{ಸಲದ} \kn{ಬಜೆಟ್} \kn{ರೈತರಿಗೆ} \kn{ಒಳ್ಳೆಯದು}.
\end{quote}

\begin{quote}
\emph{ī salada bajeṭ raitarige oḷḷeyadu.}
\end{quote}

\begin{quote}
This year's budget is good for farmers.
\end{quote}

\kn{ಮುಂಗಡ} \kn{ಪತ್ರ} is the formal Kannada name.

\begin{cousinweb}
\textbf{\kn{ಬಜೆಟ್}} (\emph{bajeṭ}) is a loan from English \emph{budget}.
\end{cousinweb}

\begin{grammarlens}[title={What you've built}]
One more word for giving reasons and opinions.
\end{grammarlens}

\subsection*{Your turn}
\begin{itemize}
\raggedright
  \item \emph{Say it:} \emph{bajeṭ}
  \item \emph{Say it:} \emph{bajeṭ}, once more
  \item \emph{Say it:} \emph{ī salada bajeṭ raitarige oḷḷeyadu.} — This year's budget is good for farmers.
  \item \emph{Recall:} say the Kannada for to distribute, then the Kannada for rural, then say \emph{bajeṭ} again
\end{itemize}

\begin{tcolorbox}[breakable,colback=teal!4,colframe=teal!35!black,title={Before you move on}]
What does \kn{ಬಜೆಟ್} mean? (\textquotedblleft{}A budget\textquotedblright{}.) Say it once more.
\end{tcolorbox}
\section[\texorpdfstring{\kn{ಆದಾಯ} (ādāya)}{ಆದಾಯ (ādāya)}]{\kn{ಆದಾಯ} (ādāya) — income}

\label{lesson:KA-C400-adaya}



\begin{quote}
Before the new one: say the Kannada for rural, then the Kannada for a budget.
\end{quote}

\subsection*{You'll want to know: \kn{ಆದಾಯ}}
\textbf{\kn{ಆದಾಯ}} — \emph{ādāya} — \textquotedblleft{}income\textquotedblright{}.

\begin{quote}
\kn{ಅವರ} \kn{ತಿಂಗಳ} \kn{ಆದಾಯ} \kn{ಕಡಿಮೆ}.
\end{quote}

\begin{quote}
\emph{avara tiṅgaḷa ādāya kaḍime.}
\end{quote}

\begin{quote}
Their monthly income is low.
\end{quote}

\kn{ಆದಾಯ} \kn{ತೆರಿಗೆ}, income tax.

\begin{cousinweb}
\textbf{\kn{ಆದಾಯ}} (\emph{ādāya}) is borrowed from Sanskrit \emph{ādāya}, \textquotedblleft{}taking, seizing\textquotedblright{} (Monier-Williams), from \emph{ā-dā}, \textquotedblleft{}to take\textquotedblright{}: what one receives.
\end{cousinweb}

\begin{grammarlens}[title={What you've built}]
One more word for giving reasons and opinions.
\end{grammarlens}

\subsection*{Your turn}
\begin{itemize}
\raggedright
  \item \emph{Say it:} \emph{ādāya}
  \item \emph{Say it:} \emph{ādāya}, once more
  \item \emph{Say it:} \emph{avara tiṅgaḷa ādāya kaḍime.} — Their monthly income is low.
  \item \emph{Recall:} say the Kannada for rural, then the Kannada for a budget, then say \emph{ādāya} again
\end{itemize}

\begin{tcolorbox}[breakable,colback=teal!4,colframe=teal!35!black,title={Before you move on}]
What does \kn{ಆದಾಯ} mean? (\textquotedblleft{}Income\textquotedblright{}.) Say it once more.
\end{tcolorbox}
\section[\texorpdfstring{\kn{ಸಬ್ಸಿಡಿ} (sabsiḍi)}{ಸಬ್ಸಿಡಿ (sabsiḍi)}]{\kn{ಸಬ್ಸಿಡಿ} (sabsiḍi) — a subsidy}

\label{lesson:KA-C400-sabsidi}



\begin{quote}
Before the new one: say the Kannada for a budget, then the Kannada for income.
\end{quote}

\subsection*{You'll want to know: \kn{ಸಬ್ಸಿಡಿ}}
\textbf{\kn{ಸಬ್ಸಿಡಿ}} — \emph{sabsiḍi} — \textquotedblleft{}a subsidy\textquotedblright{}.

\begin{quote}
\kn{ರೈತರಿಗೆ} \kn{ಗೊಬ್ಬರದ} \kn{ಮೇಲೆ} \kn{ಸಬ್ಸಿಡಿ} \kn{ಸಿಗುತ್ತದೆ}.
\end{quote}

\begin{quote}
\emph{raitarige gobbarada mēle sabsiḍi siguttade.}
\end{quote}

\begin{quote}
Farmers get a subsidy on fertiliser.
\end{quote}

The Kannada word is \kn{ಸಹಾಯಧನ}, help money.

\begin{cousinweb}
\textbf{\kn{ಸಬ್ಸಿಡಿ}} (\emph{sabsiḍi}) is a loan from English \emph{subsidy}.
\end{cousinweb}

\begin{grammarlens}[title={What you've built}]
One more word for giving reasons and opinions.
\end{grammarlens}

\subsection*{Your turn}
\begin{itemize}
\raggedright
  \item \emph{Say it:} \emph{sabsiḍi}
  \item \emph{Say it:} \emph{sabsiḍi}, once more
  \item \emph{Say it:} \emph{raitarige gobbarada mēle sabsiḍi siguttade.} — Farmers get a subsidy on fertiliser.
  \item \emph{Recall:} say the Kannada for a budget, then the Kannada for income, then say \emph{sabsiḍi} again
\end{itemize}

\begin{tcolorbox}[breakable,colback=teal!4,colframe=teal!35!black,title={Before you move on}]
What does \kn{ಸಬ್ಸಿಡಿ} mean? (\textquotedblleft{}A subsidy\textquotedblright{}.) Say it once more.
\end{tcolorbox}
\section[\texorpdfstring{\kn{ವಿನಿಯೋಗಿಸು} (viniyōgisu)}{ವಿನಿಯೋಗಿಸು (viniyōgisu)}]{\kn{ವಿನಿಯೋಗಿಸು} (viniyōgisu) — to utilise; to spend}

\label{lesson:KA-C400-viniyogisu}



\begin{quote}
Before the new one: say the Kannada for income, then the Kannada for a subsidy.
\end{quote}

\subsection*{You'll want to know: \kn{ವಿನಿಯೋಗಿಸು}}
\textbf{\kn{ವಿನಿಯೋಗಿಸು}} — \emph{viniyōgisu} — \textquotedblleft{}to utilise; to spend\textquotedblright{}.

\begin{quote}
\kn{ಪಂಚಾಯಿತಿ} \kn{ಹಣವನ್ನು} \kn{ರಸ್ತೆಗಾಗಿ} \kn{ವಿನಿಯೋಗಿಸಿತು}.
\end{quote}

\begin{quote}
\emph{pañcāyiti haṇavannu rastegāgi viniyōgisitu.}
\end{quote}

\begin{quote}
The village council spent the money on roads.
\end{quote}

Past (I …): \kn{ವಿನಿಯೋಗಿಸಿದೆ}, \emph{viniyōgiside}.

\kn{ವಿನಿಯೋಗ}, use or allocation.

\begin{cousinweb}
\textbf{\kn{ವಿನಿಯೋಗಿಸು}} (\emph{viniyōgisu}) is built on Sanskrit \emph{viniyoga}, \textquotedblleft{}apportionment, distribution, use\textquotedblright{} (Monier-Williams), with the Kannada verb ending -\emph{isu}.
\end{cousinweb}

\begin{grammarlens}[title={What you've built}]
One more word for giving reasons and opinions.
\end{grammarlens}

\subsection*{Your turn}
\begin{itemize}
\raggedright
  \item \emph{Say it:} \emph{viniyōgisu}
  \item \emph{Say it:} \emph{viniyōgisu}, once more
  \item \emph{Say it:} \emph{pañcāyiti haṇavannu rastegāgi viniyōgisitu.} — The village council spent the money on roads.
  \item \emph{Recall:} say the Kannada for income, then the Kannada for a subsidy, then say \emph{viniyōgisu} again
\end{itemize}

\begin{tcolorbox}[breakable,colback=teal!4,colframe=teal!35!black,title={Before you move on}]
What does \kn{ವಿನಿಯೋಗಿಸು} mean? (\textquotedblleft{}To utilise; to spend\textquotedblright{}.) Say it once more.
\end{tcolorbox}
\section[\texorpdfstring{\kn{ವ್ಯರ್ಥ} \kn{ಮಾಡು} (vyartha māḍu)}{ವ್ಯರ್ಥ ಮಾಡು (vyartha māḍu)}]{\kn{ವ್ಯರ್ಥ} \kn{ಮಾಡು} (vyartha māḍu) — to waste}

\label{lesson:KA-C400-vyartha-madu}



\begin{quote}
Before the new one: say the Kannada for a subsidy, then the Kannada for to utilise.
\end{quote}

\subsection*{You'll want to know: \kn{ವ್ಯರ್ಥ} \kn{ಮಾಡು}}
\textbf{\kn{ವ್ಯರ್ಥ} \kn{ಮಾಡು}} — \emph{vyartha māḍu} — \textquotedblleft{}to waste\textquotedblright{}.

\begin{quote}
\kn{ನಾನು} \kn{ಇಡೀ} \kn{ದಿನ} \kn{ಫೋನಿನಲ್ಲಿ} \kn{ಸಮಯ} \kn{ವ್ಯರ್ಥ} \kn{ಮಾಡಿದೆ}.
\end{quote}

\begin{quote}
\emph{nānu iḍī dina phōninalli samaya vyartha māḍide.}
\end{quote}

\begin{quote}
I wasted the whole day on the phone.
\end{quote}

Past (I …): \kn{ವ್ಯರ್ಥ} \kn{ಮಾಡಿದೆ}, \emph{vyartha māḍide}.

\kn{ವ್ಯರ್ಥ} alone means in vain: \kn{ವ್ಯರ್ಥವಾಯಿತು}, it was in vain.

\begin{cousinweb}
\textbf{\kn{ವ್ಯರ್ಥ} \kn{ಮಾಡು}} (\emph{vyartha māḍu}) is built from Sanskrit \emph{vyartha}, \textquotedblleft{}useless, vain\textquotedblright{} (Monier-Williams), from \emph{vi-} \textquotedblleft{}without\textquotedblright{} and \emph{artha} \textquotedblleft{}purpose\textquotedblright{}, with \emph{māḍu}, \textquotedblleft{}to do\textquotedblright{}.
\end{cousinweb}

\begin{grammarlens}[title={What you've built}]
One more word for giving reasons and opinions.
\end{grammarlens}

\subsection*{Your turn}
\begin{itemize}
\raggedright
  \item \emph{Say it:} \emph{vyartha māḍu}
  \item \emph{Say it:} \emph{vyartha māḍu}, once more
  \item \emph{Say it:} \emph{nānu iḍī dina phōninalli samaya vyartha māḍide.} — I wasted the whole day on the phone.
  \item \emph{Recall:} say the Kannada for a subsidy, then the Kannada for to utilise, then say \emph{vyartha māḍu} again
\end{itemize}

\begin{tcolorbox}[breakable,colback=teal!4,colframe=teal!35!black,title={Before you move on}]
What does \kn{ವ್ಯರ್ಥ} \kn{ಮಾಡು} mean? (\textquotedblleft{}To waste\textquotedblright{}.) Say it once more.
\end{tcolorbox}
